% CAIS 2026 Workshop Extended Abstract
% ACM Conference on AI Safety, May 27-29, San Jose
% Draft v7: April 10, 2026
% V7: Integrate Claude Chorus LLM experiment, remove bridge (for evolving-strategies),
%   add contravariant functor explanation, explicit agent definitions.
%   Addresses Robin's UID 975/977 feedback: agent-side quantities defined explicitly.
\documentclass[10pt,twocolumn]{article}

\usepackage[margin=1in]{geometry}
\usepackage{graphicx}
\usepackage{booktabs}
\usepackage{amsmath,amssymb}
\usepackage{amsthm}
\usepackage[T1]{fontenc}
\usepackage[hyphens]{url}
\usepackage[expansion=false]{microtype}
\usepackage{natbib}
\usepackage{hyperref}
\usepackage{enumitem}
\setlist{nosep,leftmargin=*}

\newtheorem{definition}{Definition}

\graphicspath{{figures/}}

\title{What Your Agent Topology Is Actually Doing:\\
Controlled Experiments on Communication Structure\\
in Multi-Agent Systems}

\author{
  Robin Langer \and Lyra Vega
}

\date{}

\begin{document}
\maketitle

\begin{abstract}
Multi-agent AI systems choose communication topologies---star, chain,
ring, fully connected---with no controlled evidence for how topology
affects behavior. We provide three lines of evidence. First, using
island-model genetic algorithms as a model system (240~runs), we show
topology independently explains 17\% of behavioral diversity variance
($p < 10^{-7}$), amplified $14\times$ on hard problems. Second, we
validate the transfer to real agents: 1,600 LLM~API calls across
8~coding tasks show that topology determines solution diversity in
GPT-4o-mini agents with \emph{inverted} sign---isolated agents
converge (Cohen's $d = -11.1$), while fully connected agents diversify
(Kendall's $W = 0.91$). The inversion arises because GA migration
homogenizes (left Kan extension) while LLM context-sharing diversifies
(right Kan extension)---same topology, opposite effect, determined by
which adjoint the coupling mechanism instantiates. At least nine
deployed agent systems chose topologies based on uncontrolled
comparisons. Topology is a first-class design variable for agent
safety; current systems are tuning it blindly.
\end{abstract}

%% ===================================================================
\section{Introduction}
\label{sec:intro}
%% ===================================================================

Every multi-agent system embeds a communication topology---round-robin
group chat (AutoGen), DAG pipelines (LangGraph), manager hubs
(CrewAI)---yet these choices are made without controlled evidence.
Dochkina~\citep{dochkina2026} ran 25,000 tasks across 256 agents,
finding a 44\% quality spread ($d = 1.86$) between topologies.
Hierarchical deployments fail 36\% of the time; swarms fail 68\%.
Topology is doing something---but \emph{what, exactly?}

We provide the first controlled answer, bridging from a model system
(genetic algorithms) to real LLM agents. Our contributions:
(1)~topology independently determines behavioral diversity in both GAs
and LLM agents; (2)~the effect \emph{inverts} between GAs and
agents, explained by which adjoint of the topology the coupling
mechanism instantiates; (3)~at least nine deployed systems chose
topologies without controlled comparisons.

%% ===================================================================
\section{Island Models as Agent Systems}
\label{sec:analogy}
%% ===================================================================

An island-model GA partitions a population across sub-populations
(\emph{islands}) that evolve independently, exchanging individuals
periodically via \emph{migration} channels. The mapping to agent
orchestration is direct:

\begin{center}
\small
\begin{tabular}{@{}p{3.2cm}p{3.8cm}@{}}
\toprule
\textbf{GA concept} & \textbf{Agent equivalent} \\
\midrule
Island & Agent (local state, computation) \\
Migration channel & Message channel \\
Population diversity & Solution diversity \\
Fitness function & Task performance metric \\
Migration topology & Orchestration graph \\
\bottomrule
\end{tabular}
\end{center}

\noindent
The analogy is structural, not metaphorical: both systems have nodes
with local state exchanging information along directed channels.
The agent-side quantities are measured directly:
\textbf{agent diversity} = mean pairwise cosine distance on sentence
embeddings (all-MiniLM-L6-v2) of agent outputs;
\textbf{agent fitness} = task performance (test pass rate,
correctness). GAs provide cheap experimental control (hundreds of
trials); the LLM experiment (Section~\ref{sec:chorus}) validates
transfer to real agents.

%% ===================================================================
\section{Definitions}
\label{sec:definitions}
%% ===================================================================

\begin{definition}[Communication topology]
A \emph{communication topology} for $n$ agents is a directed graph
$D = (V, A)$ with $|V| = n$, where each arc $(u, v) \in A$ is a
permitted channel from agent $u$ to agent $v$.
\end{definition}

\begin{definition}[Cycle rank]
The \emph{cycle rank} $\beta_1(G) = |E| - |V| + 1$ (for connected
$G$) counts independent cycles---feedback loops in the communication
structure.
\end{definition}

\begin{definition}[Directed cycle count]
$\kappa(D)$ counts distinct simple directed cycles in
$D$~\citep{johnson1975finding}. Unlike $\beta_1$, this depends on arc
\emph{arrangement}, not just edge count. A DAG has $\kappa = 0$.
\end{definition}

\paragraph{NK landscapes}~\citep{kauffman1993origins} are tunably
rugged fitness functions over binary strings of length $N$, where each
bit's contribution depends on $K$ other bits. $K{=}0$ gives a smooth
landscape; increasing $K$ adds local optima.

%% ===================================================================
\section{Experimental Evidence}
\label{sec:experiments}
%% ===================================================================

We present three experiments spanning GA model systems and LLM agents,
totaling 2,400+ runs and API calls.

\subsection{Experiment 1: Topology affects diversity at constant density}

Comparing a chain to a complete graph conflates density and cycle
structure: $\beta_1 = m - n + 1$~\citep{hatcher2002algebraic} means
cycle count and edge count are perfectly correlated for undirected
graphs. We escape via directed cycle count $\kappa$, which varies
independently of density.

\paragraph{Design.}
8 directed graphs, all $n{=}8$, 16 arcs ($\delta{=}2.0$), varying
only in $\kappa \in \{0, 0, 3, 10, 14, 20, 29, 47\}$.
Island-model GA: 50 individuals/island, 500 generations, 30 seeds
per topology, 240 runs total.

\paragraph{Results.}
$F(7, 232) = 6.90$, $p = 1.85 \times 10^{-7}$, $\eta^2 = 0.17$.
More directed cycles \emph{accelerate consensus} ($r = -0.68$):
topologies with more feedback loops cause agents to converge faster,
reducing capacity to maintain alternative strategies.

\subsection{Experiment 2: Hard problems amplify the effect $14\times$}

On NK landscapes with tunable ruggedness, topology's effect scales
with difficulty:

\begin{center}
\small
\begin{tabular}{@{}ccc@{}}
\toprule
$K$ & $\eta^2$ (diversity) & Interpretation \\
\midrule
0 (trivial) & 0.05 & Topology barely matters \\
2 (moderate) & 0.20 & Moderate effect \\
4 (hard) & 0.45 & Large effect \\
6 (very hard) & 0.69 & Topology dominates \\
\bottomrule
\end{tabular}
\end{center}

\noindent
Simple tasks are structurally smooth---topology barely matters.
Complex tasks (multiple valid strategies, ambiguous evaluation) are
rugged---topology dominates. Real-world deployments overwhelmingly
involve the latter.

\subsection{Experiment 3: Topology inverts diversity in LLM agents}
\label{sec:chorus}

The GA experiments establish that topology controls diversity.
But does this transfer to LLM agents? We test directly.

\paragraph{Design.}
5 GPT-4o-mini agents, 8 coding tasks (shortest path, kth largest,
rate limiter, JSON parser, LRU cache, regex matcher, expression
evaluator, Bloom filter), 4 topologies (none, ring, star, fully
connected), 10 seeds per condition. 1,600 API calls total.
Coupling is \emph{negative}: each agent is instructed to use a
\emph{different} algorithm from those visible in its context. Diversity
is measured as mean pairwise cosine distance on sentence embeddings
(all-MiniLM-L6-v2) of agent outputs.

\paragraph{Results.}
The diversity ordering \textbf{inverts} from the GA result:

\begin{center}
\small
\begin{tabular}{@{}llp{3.6cm}@{}}
\toprule
\textbf{System} & \textbf{Diversity ordering} & \textbf{Statistic} \\
\midrule
GA & none $>$ ring $>$ FC & $W = 1.0$ (6 domains) \\
LLM & FC $>$ ring $\approx$ star $>$ none & $W = 0.91$ (6/8 tasks) \\
\bottomrule
\end{tabular}
\end{center}

\noindent
Six of eight tasks show the inverted ordering cleanly (Kruskal-Wallis
$p < 10^{-6}$, Cohen's $d = -11.1$ for none vs.\ fully connected on
the strongest task). Per-task results:

\begin{center}
\footnotesize
\begin{tabular}{@{}lccccr@{}}
\toprule
\textbf{Task} & \textbf{None} & \textbf{Ring} & \textbf{Star} & \textbf{FC} & $d$ \\
\midrule
shortest\_path & .041 & .174 & .116 & \textbf{.231} & $-$14.8 \\
kth\_largest & .043 & .199 & .157 & \textbf{.244} & $-$8.8 \\
rate\_limiter & .029 & .096 & .093 & \textbf{.133} & $-$5.9 \\
json\_parser & .088 & .126 & .137 & \textbf{.149} & $-$3.5 \\
lru\_cache & .071 & .127 & .137 & \textbf{.172} & $-$2.6 \\
regex\_match & .143 & .183 & .150 & \textbf{.205} & $-$1.9 \\
\midrule
expr\_eval & .254 & .289 & .278 & .232 & $+$0.4 \\
bloom\_filter & .047 & .052 & .046 & .050 & $\approx$0 \\
\bottomrule
\end{tabular}
\end{center}

\noindent
The two exceptions are informative. \textbf{Bloom filter} (spread
$= 0.006$): the task is unimodal---no algorithmic variety for
topology to act on. \textbf{Expression evaluator} ($d = +0.4$):
the model's prior is already spread across multiple approaches, so
the ``be different'' constraint adds little.

\paragraph{Why the sign flips.}
Under ``none'' topology, all 5 LLM agents independently generate
the \emph{same} algorithm (e.g., quickselect for kth\_largest).
LLMs start from a \textbf{peaked prior} (low entropy); GAs start
from \textbf{random initialization} (high entropy). Isolation
preserves initial entropy: high for GAs (diverse), low for LLMs
(homogeneous). The inversion follows from which \emph{adjoint} of the
topology the coupling mechanism instantiates:

\begin{itemize}
  \item \textbf{GA migration} = \textbf{left Kan extension}
    (colimit-like): merges populations. More connectivity $\to$ less
    diversity.
  \item \textbf{LLM ``be different'' context} = \textbf{right Kan
    extension} (limit-like): accumulates exclusion constraints. More
    connectivity $\to$ agents forced into distributional tails $\to$
    \emph{more} diversity.
\end{itemize}

\noindent
The topology is the invariant; the sign depends on the adjoint.
$\mathrm{Lan}_F \dashv \mathrm{Ran}_F$ reverses ordering on any poset,
including the diversity ordering on topologies.

%% ===================================================================
\section{Implications for Agent Orchestration}
\label{sec:implications}
%% ===================================================================

\paragraph{The evidence gap.}
At least nine recent agent systems---spanning ICML, AAAI, WWW, and
NeurIPS---chose topologies without density-controlled comparisons
(Table~\ref{tab:systems}).

\begin{table}[t]
\centering
\caption{Agent systems whose topology choices lack controlled evidence.}
\label{tab:systems}
\footnotesize
\resizebox{\columnwidth}{!}{%
\begin{tabular}{@{}lp{3.5cm}@{}}
\toprule
\textbf{System} & \textbf{Issue} \\
\midrule
Graph-GRPO~\citep{graphgrpo2026} & Complete = worst; also max $\beta_1$ \\
OFA-MAS~\citep{ofamas2026} & Prunes density \& cycles together \\
G-Designer~\citep{gdesigner2025} & No density-controlled baseline \\
GoAgent~\citep{goagent2026} & DAG vs.\ dense alternatives \\
ARG-Designer~\citep{argdesigner2026} & DAG space only \\
AgentConductor~\citep{agentconductor2026} & Cycles never tested \\
DyTopo~\citep{dytopo2026} & DAG enforced throughout \\
HERA~\citep{hera2026} & Cycles tracked, not controlled \\
Dochkina~\citep{dochkina2026} & Density \& cycles confounded \\
\bottomrule
\end{tabular}}
\end{table}

\paragraph{Safety risks.}
Our LLM experiment sharpens three concerns:
\emph{(1)~The sign matters.} A topology that reduces diversity in
GAs may \emph{increase} it in agents with negative coupling (``be
different'' instructions). Predictions based on GA intuition alone
will have the wrong sign.
\emph{(2)~Premature consensus.} Isolated LLM agents converge to
identical solutions (Cohen's $d = -11.1$ vs.\ connected agents).
``Ensemble diversity'' claims based on independent agents may be
illusory when models share a peaked prior.
\emph{(3)~Error amplification.} Moran~\citep{moran2025error} reports
$17\times$ error amplification in unstructured communication. The
standard fix (impose DAG structure) simultaneously removes cycles
\emph{and} density---we cannot know which helps.

\paragraph{Recommendations.}
\begin{enumerate}
  \item \textbf{Treat topology as a first-class variable.}
    It controls the diversity--convergence tradeoff---in both GAs and
    LLM agents, empirically confirmed.
  \item \textbf{Know your coupling sign.} Does your agent framework
    encourage agents to build on neighbors' outputs (homogenizing,
    left Kan) or differentiate from them (diversifying, right Kan)?
    The same topology has opposite effects.
  \item \textbf{Match topology to task complexity.} Simple tasks
    tolerate any topology; hard tasks demand deliberate choice
    ($\eta^2$ scales $14\times$ from $K{=}0$ to $K{=}6$).
  \item \textbf{Monitor diversity, not just accuracy.} Topology
    effects manifest in strategy diversity. Accuracy alone masks
    topology-induced failures.
\end{enumerate}

%% ===================================================================
\section{Related Work}
\label{sec:related}
%% ===================================================================

Yang et al.~\citep{yang2025topology} identify topology as the
underexplored MAS dimension; we provide the controlled experiments
they call for. HERA~\citep{hera2026} tracks cycle count during
topology evolution, converging on ``compact chains with preserved
cycles.'' Dochkina~\citep{dochkina2026} shows massive effects (25K
tasks, $d = 1.86$) without isolating the mechanism. In evolutionary
computation, the topology--diversity link is
established~\citep{tomassini2005spatially,skolicki2005influence} but
without density control. Capucci and
Myers~\citep{capucci2026feedback} show each cycle requires an
independent convergence guarantee---the theoretical cost side.

%% ===================================================================
\section{Conclusion}
\label{sec:conclusion}
%% ===================================================================

Communication topology is not neutral infrastructure---it is a design
variable controlling how multi-agent systems balance exploration
against convergence. In GA model systems, topology independently
explains 17\% of diversity variance on easy problems, rising to 69\%
on hard ones. In LLM agents, the effect is equally strong but
\emph{inverted}: fully connected agents produce the most diverse
solutions (Cohen's $d = -11.1$), while isolated agents converge to
identical outputs. The sign depends on which adjoint of the topology
graph the coupling mechanism instantiates---left Kan (merge,
homogenize) vs.\ right Kan (constrain, diversify). At least nine
deployed systems chose topologies based on comparisons that
simultaneously varied the property they claimed to test. For safe
multi-agent deployment, topology must be tuned deliberately, its
coupling sign understood, and its effects monitored through diversity
metrics---not assumed away.

\smallskip
\noindent\textbf{Acknowledgments.} [Withheld for review.]

%% ===================================================================
%% References
%% ===================================================================

\bibliographystyle{plainnat}
\begin{thebibliography}{20}

\bibitem[Capucci and Myers(2026)]{capucci2026feedback}
M.~Capucci and D.~J. Myers.
\newblock Feedback in categorical multi-agent systems.
\newblock \emph{arXiv:2604.03017}, 2026.

\bibitem[Dochkina(2026)]{dochkina2026}
V.~Dochkina.
\newblock Drop the hierarchy and roles: How self-organizing {LLM} agents
  outperform designed structures.
\newblock \emph{arXiv:2603.28990}, 2026.

\bibitem[Chen et~al.(2026)]{goagent2026}
H.~Chen et~al.
\newblock GoAgent: Group-of-agents communication topology generation for
  {LLM}-based multi-agent systems.
\newblock \emph{arXiv:2603.19677}, 2026.

\bibitem[Wang et~al.(2026)]{agentconductor2026}
S.~Wang et~al.
\newblock AgentConductor: Topology evolution for multi-agent competition-level
  code generation.
\newblock \emph{arXiv:2602.17100}, 2026.

\bibitem[Cang et~al.(2026)]{graphgrpo2026}
Y.~Cang et~al.
\newblock Graph-GRPO: Stabilizing multi-agent topology learning via group
  relative policy optimization.
\newblock \emph{arXiv:2603.02701}, 2026.

\bibitem[Li et~al.(2026)]{ofamas2026}
S.~Li et~al.
\newblock OFA-MAS: One-for-all multi-agent system topology design based on
  mixture-of-experts graph generative models.
\newblock In \emph{Proc.\ WWW}, 2026.

\bibitem[Zhang et~al.(2025)]{gdesigner2025}
G.~Zhang et~al.
\newblock G-Designer: Architecting multi-agent communication topologies via
  graph neural networks.
\newblock In \emph{Proc.\ ICML}, 2025.

\bibitem[Li et~al.(2026b)]{argdesigner2026}
S.~Li et~al.
\newblock Assemble your crew: Automatic multi-agent communication topology
  design via autoregressive graph generation.
\newblock In \emph{Proc.\ AAAI}, 2026.

\bibitem[Hatcher(2002)]{hatcher2002algebraic}
A.~Hatcher.
\newblock \emph{Algebraic Topology}.
\newblock Cambridge University Press, 2002.

\bibitem[Li and Ramakrishnan(2026)]{hera2026}
S.~Li and N.~Ramakrishnan.
\newblock Experience as a compass: Multi-agent {RAG} with evolving
  orchestration and agent prompts.
\newblock \emph{arXiv:2604.00901}, 2026.

\bibitem[Johnson(1975)]{johnson1975finding}
D.~B. Johnson.
\newblock Finding all the elementary circuits of a directed graph.
\newblock \emph{SIAM J.\ Comput.}, 4(1):77--84, 1975.

\bibitem[Kauffman(1993)]{kauffman1993origins}
S.~A. Kauffman.
\newblock \emph{The Origins of Order}.
\newblock Oxford University Press, 1993.

\bibitem[Moran(2025)]{moran2025error}
K.~Moran.
\newblock The secret is topology: Why unstructured multi-agent communication
  amplifies errors.
\newblock \emph{Towards Data Science}, 2025.

\bibitem[Skolicki and De~Jong(2005)]{skolicki2005influence}
Z.~Skolicki and K.~A. De~Jong.
\newblock The influence of migration sizes and intervals on island models.
\newblock In \emph{GECCO}, pages 1295--1302, 2005.

\bibitem[Tomassini(2005)]{tomassini2005spatially}
M.~Tomassini.
\newblock \emph{Spatially Structured Evolutionary Algorithms}.
\newblock Springer, 2005.

\bibitem[Yang et~al.(2025)]{yang2025topology}
J.~Yang et~al.
\newblock Topological structure learning should be a research priority for
  {LLM}-based multi-agent systems.
\newblock \emph{arXiv:2505.22467}, 2025.

\bibitem[Lu et~al.(2026)]{dytopo2026}
Y.~Lu et~al.
\newblock DyTopo: Dynamic topology routing for multi-agent reasoning via
  semantic matching.
\newblock \emph{arXiv:2602.06039}, 2026.

\end{thebibliography}

\end{document}
